% Danica's UBC CV
% Adapted from Kwang Moo Yi's and Kevin-Leyton-Brown's CVs,
% which had influences from Ron Garcia's and Helge Rhodin's

\documentclass[10pt]{article}


% set up toggles for format that we use early
\usepackage{etoolbox}
\newtoggle{ubcformat}
\newtoggle{annualReview}

% set them based on jobname (e.g. latexmk -jobname=cv ubc-cv.tex)
% https://texfaq.org/FAQ-compjobnam
\def\StripPrefix#1>{}
\def\jobis#1{FF\fi
  \def\predicate{#1}%
  \edef\predicate{\expandafter\StripPrefix\meaning\predicate}%
  \edef\job{\jobname}%
  \ifx\job\predicate
}

\if\jobis{ubc-cv}%
  \toggletrue{ubcformat}%
\else
    \if\jobis{ubc-cv-anno}%
        \toggletrue{ubcformat}%
        \toggletrue{annualReview}%
    \else
    \fi
\fi


% Adapted from Ron Garcia's UBC CV
% Adapted from Kevin Leyton-Brown's CV Template
% Adapted from Helge Rhodin's CV Template

\usepackage[USenglish]{babel}
\usepackage[utf8]{inputenc}
\usepackage{csquotes}

\usepackage{longtable}  % Support tables that span pages
\usepackage{xcolor}
\usepackage{hhline}
\usepackage{booktabs} % Commands for book-quality tables
\usepackage{comment} % Enable and disable "blurbs" (for Evaluative CVs)
\usepackage{multirow} % Better looking multi-row table headers
\usepackage{enumitem} % Better list control
\usepackage{lipsum} % Lipsum

% Font to arial
\usepackage{helvet}
\renewcommand{\familydefault}{\sfdefault}

% Cool underline
% https://tex.stackexchange.com/questions/443439/underline-width
\usepackage{soul}
\newcommand\dunderline[3][-1pt]{
  {%
      \sbox0{#3}%
      \ooalign{\copy0\cr\rule[\dimexpr#1-#2\relax]{\wd0}{#2}}
    }
}

%%
%% Logo / Header
%%
\usepackage{graphicx}
\newcommand{\UBClogo}{\includegraphics[height=9pt]{figs/UBC-logo-grey}~}
\newcommand{\bigUBClogo}{\includegraphics[height=26pt]{figs/UBC-logo-grey}~}
\newcommand{\UBCheader}[1]{\begin{center} \begin{minipage}{.065\textwidth}\bigUBClogo\end{minipage}~\begin{minipage}{.5\textwidth}\centering{\includegraphics[width=\textwidth]{figs/ubc-logo-2018-wordmark-black-cmyk}}\\{\bf\it #1}\end{minipage}~~\begin{minipage}{.1\textwidth}~\end{minipage}\end{center}}
%\newcommand{\UBCheader}[1]{\begin{center} \begin{minipage}{.1\textwidth}\bigUBClogo\end{minipage}~~\begin{minipage}{.5\textwidth}\centering{\bf THE UNIVERSITY OF BRITISH COLUMBIA} \\{\bf\it #1}\end{minipage}~~\begin{minipage}{.1\textwidth}~\end{minipage}\end{center}}


% here's how to use fixed footnotes with symbol fonts!  https://tex.stackexchange.com/questions/68703/fixfoot-sty-with-symbols
\usepackage[symbol*,perpage,bottom]{footmisc}
\usepackage{fixfoot}
\usepackage{etoolbox}% provides \patchcmd
\makeatletter
\patchcmd\@fixed@footnote
{\protected@xdef\@thefnmark{\csname @#1@fftn@footnote\endcsname}}% search
{\protected@xdef\@thefnmark{%
		\expandafter\@fnsymbol\csname @#1@fftn@footnote\endcsname}}% replace
{}{}% success/failure
\makeatother

\DeclareFixedFootnote{\paid}{Expenses paid.}
\DeclareFixedFootnote{\competitive}{Competitively selected.}



%%
%% Annual Review related annotations
%%

% Comment below to disable annual review annotations
\usepackage{colortbl}
\usepackage{marginnote}
\usepackage{etoolbox}
\newtoggle{annualReview}
\toggletrue{annualReview}
\iftoggle{annualReview}
{
  \definecolor{newcolor}{rgb}{1,0,0}
  \definecolor{newtablecolor}{rgb}{1,.627,.627}
  \definecolor{continuingcolor}{rgb}{0,.353,1}
  \definecolor{continuingtablecolor}{rgb}{.612,.733,1}
  %    \definecolor{changedcolor}{rgb}{.376,.290,.482}
  %    \definecolor{changedtablecolor}{rgb}{.702,.635,.780}
  \definecolor{toappearcolor}{rgb}{0,.353,1}
  \definecolor{toappeartablecolor}{rgb}{.612,.733,1}
  \definecolor{changedcolor}{RGB}{18,135,25}
  \definecolor{changedtablecolor}{RGB}{141,202,145}
  \definecolor{acceptedcolor}{RGB}{18,135,25}
  \definecolor{acceptedtablecolor}{RGB}{141,202,145}
  \definecolor{boxcolor}{rgb}{.9,.9,.9}
  \newcommand{\newpic}{\includegraphics[height=9pt]{figs/new}}
  \newcommand{\acceptedpic}{\includegraphics[height=9pt]{figs/accepted}}
  \newcommand{\toappearpic}{\includegraphics[height=9pt]{figs/toappear}}
  \newcommand{\continuingpic}{\includegraphics[height=9pt]{figs/continuing}}
  \newcommand{\changedpic}{\includegraphics[height=9pt]{figs/changed}}
  \newcommand{\new}[1]{\marginnote{\newpic}{\color{newcolor} #1}}
  \newcommand{\accepted}[1]{\marginnote{\acceptedpic}{\color{acceptedcolor} #1}}
  \newcommand{\toappear}[1]{\marginnote{\toappearpic}{\color{toappearcolor} #1}}
  \newcommand{\continuing}[1]{\marginnote{\continuingpic}{\color{continuingcolor} #1}}
  \newcommand{\changed}[1]{\marginnote{\changedpic}{\color{changedcolor} #1}}
  \newcommand{\newtable}{\rowcolor{newtablecolor}\marginnote{\newpic}}
  \newcommand{\continuingtable}{\rowcolor{continuingtablecolor}\marginnote{\continuingpic}}
  \newcommand{\changedtable}{\rowcolor{changedtablecolor}\marginnote{\changedpic}}
  \newcommand{\toonew}[1]{}

  \reversemarginpar
  \usepackage[letterpaper,left=2.5cm,right=2cm,top=.5in,bottom=.5in,marginparwidth=1.7cm,includeheadfoot,headsep=16pt]{geometry}
}
{
  \newcommand{\new}[1]{#1}
  \newcommand{\changed}[1]{#1}
  \newcommand{\continuing}[1]{#1}
  \newcommand{\newtable}{}
  \newcommand{\changedtable}{}
  \newcommand{\continuingtable}{}

  \usepackage[letterpaper,left=2.25cm,right=2.25cm,top=.5in,bottom=.5in,marginparwidth=0cm,includeheadfoot,headsep=16pt]{geometry}
}


%%
%% Tables
%%
\usepackage{tabularx}
\usepackage{ragged2e}

\newenvironment{coursetable}{
  \begin{center}
    \begin{longtable}[h!]{ p{0.1\textwidth} p{0.15\textwidth} p{0.07\textwidth} p{0.07\textwidth} p{0.1\textwidth} p{0.1\textwidth} p{0.09\textwidth} p{0.19\textwidth} }
      %% The header for the top of the table
      \toprule
      \multicolumn{1}{c}{\bf Session} & \multicolumn{1}{c}{\bf Course} & \multicolumn{1}{c}{\bf Sched.} & \multicolumn{1}{c}{\bf Class} & \multicolumn{4}{c}{\bf Hours Taught}                                                                                                    \\
      \cmidrule(lr){5-8}
                                      & \multicolumn{1}{c}{\bf Number} & \multicolumn{1}{c}{\bf Hours}  & \multicolumn{1}{c}{\bf Size}  & \multicolumn{1}{c}{\bf Lectures}     & \multicolumn{1}{c}{\bf Tutorials} & \multicolumn{1}{c}{\bf Labs} & \multicolumn{1}{c}{\bf Other} \\
      \midrule \endfirsthead
      %% The header for the continuation pages
      \toprule
      \multicolumn{1}{c}{\bf Session} & \multicolumn{1}{c}{\bf Course} & \multicolumn{1}{c}{\bf Sched.} & \multicolumn{1}{c}{\bf Class} & \multicolumn{4}{c}{\bf Hours Taught}                                                                                                    \\
      \cmidrule(lr){5-8}
                                      & \multicolumn{1}{c}{\bf Number} & \multicolumn{1}{c}{\bf Hours}  & \multicolumn{1}{c}{\bf Size}  & \multicolumn{1}{c}{\bf Lectures}     & \multicolumn{1}{c}{\bf Tutorials} & \multicolumn{1}{c}{\bf Labs} & \multicolumn{1}{c}{\bf Other} \\
      \midrule \endhead
      %% Footer for the bottom of the table
      \bottomrule \endlastfoot
      %% Footer for the continued pages
      \bottomrule \footnotesize{\emph{(continued\ldots)}} \endfoot
      }{
    \end{longtable}
  \end{center}
}

\newenvironment{suptable}{
  \begin{center}
    \begin{longtable}[h!]{ p{0.26\textwidth} p{0.25\textwidth} p{0.07\textwidth} p{0.07\textwidth} p{0.27\textwidth} }
      %% The header for the top of the table
      \toprule
      \multicolumn{1}{c}{\bf Student} & \multicolumn{1}{c}{\bf Program} & \multicolumn{1}{c}{\bf Start} & \multicolumn{1}{c}{\bf Finish}  & \multicolumn{1}{c}{\bf Supervisory Role}                                                                                            \\
      % \cmidrule(lr){3-4}
      %                                      & \multicolumn{1}{c}{\bf Type}    & \multicolumn{1}{c}{\bf Start} & \multicolumn{1}{c}{\bf Finish}           & \multicolumn{1}{c}{\bf\parbox[t]{0.27\textwidth}{\centering(supervisor, co-supervisor)}} \\
      \midrule \endfirsthead
      %% The header for the continuation pages
      \toprule
      \multicolumn{1}{c}{\bf Student} & \multicolumn{1}{c}{\bf Program} & \multicolumn{1}{c}{\bf Start} & \multicolumn{1}{c}{\bf Finish}  & \multicolumn{1}{c}{\bf Supervisory Role}                                                                                            \\
      % \cmidrule(lr){3-4}
      %                                      & \multicolumn{1}{c}{\bf Type}    & \multicolumn{1}{c}{\bf Start} & \multicolumn{1}{c}{\bf Finish}           & \multicolumn{1}{c}{\bf\parbox[t]{0.27\textwidth}{\centering(supervisor, co-supervisor)}} \\
      \midrule \endhead
      %% Footer for the bottom of the table
      \bottomrule \endlastfoot
      %% Footer for the continued pages
      \bottomrule \footnotesize{\emph{(continued\ldots)}} \endfoot
      }{
    \end{longtable}
  \end{center}
}


\newenvironment{supcommtable}{
  \begin{center}
    \begin{longtable}[h!]{ p{0.23\textwidth} p{0.20\textwidth} p{0.07\textwidth} p{0.07\textwidth} p{0.3\textwidth} }
      %% The header for the top of the table
      \toprule
      \multicolumn{1}{c}{\bf Student Name} & \multicolumn{1}{c}{\bf Program} & \multicolumn{2}{c}{\bf Year}  & \multicolumn{1}{c}{\bf Committee Type}                                                                                \\
      \cmidrule(lr){3-4}
                                           & \multicolumn{1}{c}{\bf Type}    & \multicolumn{1}{c}{\bf Start} & \multicolumn{1}{c}{\bf Finish}         & \multicolumn{1}{c}{\bf\parbox[t]{0.26\textwidth}{\centering(RPE, PhD, MSc)}} \\
      \midrule \endfirsthead
      %% The header for the continuation pages
      \toprule
      \multicolumn{1}{c}{\bf Student Name} & \multicolumn{1}{c}{\bf Program} & \multicolumn{2}{c}{\bf Year}  & \multicolumn{1}{c}{\bf Committee Type}                                                                                \\
      \cmidrule(lr){3-4}
                                           & \multicolumn{1}{c}{\bf Type}    & \multicolumn{1}{c}{\bf Start} & \multicolumn{1}{c}{\bf Finish}         & \multicolumn{1}{c}{\bf\parbox[t]{0.26\textwidth}{\centering(RPE, PhD, MSc)}} \\
      \midrule \endhead
      %% Footer for the bottom of the table
      \bottomrule \endlastfoot
      %% Footer for the continued pages
      \bottomrule \footnotesize{\emph{(continued\ldots)}} \endfoot
      }{
    \end{longtable}
  \end{center}
}

\newenvironment{lecttable}{
  \begin{center}
    \begin{longtable}[h!]{ p{0.1\textwidth} p{0.2\textwidth} p{0.07\textwidth} p{0.14\textwidth} p{0.375\textwidth} }
      %% The header for the top of the table
      \toprule
      \multicolumn{1}{c}{\bf Date} & \multicolumn{1}{c}{\bf Event Name} & \multicolumn{1}{c}{\bf Hours} & \multicolumn{1}{c}{\bf Audience} & \multicolumn{1}{c}{\bf Title} \\
      \midrule \endfirsthead
      %% The header for the continuation pages
      \toprule
      \multicolumn{1}{c}{\bf Date} & \multicolumn{1}{c}{\bf Event Name} & \multicolumn{1}{c}{\bf Hours} & \multicolumn{1}{c}{\bf Audience} & \multicolumn{1}{c}{\bf Title} \\
      \midrule \endhead
      %% Footer for the bottom of the table
      \bottomrule \endlastfoot
      %% Footer for the continued pages
      \bottomrule \footnotesize{\emph{(continued\ldots)}} \endfoot
      }{
    \end{longtable}
  \end{center}
}

\newenvironment{granttable}{
  \begin{center}
    \small
    \begin{longtable}[h!]{
        p{0.165\textwidth}
        p{0.24\textwidth}
        >{\RaggedLeft\arraybackslash}p{0.07\textwidth}
        >{\RaggedLeft\arraybackslash}p{0.08\textwidth}
        p{0.075\textwidth}
        >{\RaggedRight\arraybackslash}p{0.24\textwidth}
      }
      %% The header for the top of the table
      \toprule
      \textbf{Agency, Program} &
      \textbf{Title} &
      \textbf{Comp} &
      \textbf{\$/year} &
      \textbf{Year(s)} &
      \textbf{PI \emph{\small{(Co-PIs in italics)}}} \\
      \midrule \endfirsthead
      %% The header for the continuation pages
      \toprule
      \textbf{Agency, Program} &
      \textbf{Title} &
      \textbf{Comp} &
      \textbf{\$/year} &
      \textbf{Year(s)} &
      \textbf{PI \emph{\small{(Co-PIs in italics)}}} \\
      \midrule \endhead
      %% Footer for the bottom of the table
      \bottomrule \endlastfoot
      %% Footer for the continued pages
      \bottomrule \footnotesize{\emph{(continued\ldots)}} \endfoot
      }{
    \end{longtable}
  \end{center}
}

\newenvironment{extcomtable}{
  \begin{center}
    \small
    \begin{longtable}[h!]{
        p{0.05\textwidth}
        p{0.2\textwidth}
        p{0.25\textwidth}
        p{0.35\textwidth}
      }
      %% The header for the top of the table
      \toprule
      \textbf{Year} &
      \textbf{Student Name} &
      \textbf{Institution} &
      \textbf{Degree} \\
      \midrule \endfirsthead
      %% The header for the continuation pages
      \toprule
      \textbf{Year} &
      \textbf{Student Name} &
      \textbf{Institution} &
      \textbf{Degree} \\
      \midrule \endhead
      %% Footer for the bottom of the table
      \bottomrule \endlastfoot
      %% Footer for the continued pages
      \bottomrule \footnotesize{\emph{(continued\ldots)}} \endfoot
      }{
    \end{longtable}
  % \begin{center}
  %   \begin{longtable}[h!]{
  %       p{0.2\textwidth}
  %       p{0.05\textwidth}
  %       p{0.2\textwidth}
  %       p{0.2\textwidth}
  %     }
  %     %% The header for the top of the table
  %     \toprule
  %     {\bf Student Name}
  %     & {\bf Year}
  %     & {\bf Institution}
  %     & {\bf Degree}
  %     \\
  %     \midrule \endfirsthead
  %     %% The header for the continuation pages
  %     \toprule
  %     {\bf Student Name}
  %     & {\bf Year}
  %     & {\bf Institution}
  %     & {\bf Degree}
  %     \midrule \endhead
  %     %% Footer for the bottom of the table
  %     \bottomrule \endlastfoot
  %     %% Footer for the continued pages
  %     \bottomrule \footnotesize{\emph{(continued\ldots)}} \endfoot
  %     }{
  %   \end{longtable}
  \end{center}
}



% https://tex.stackexchange.com/questions/5017/center-column-with-specifying-width-in-table-tabular-enviroment
% \newcolumntype{x}[1]{>{\centering\arraybackslash\hspace{0pt}}p{#1\textwidth}}
\newcolumntype{x}[1]{>{\centering\let\newline\\\arraybackslash\hspace{0pt}}p{#1}}

% Dashed line
\usepackage{arydshln}



% Better looking two-line headers. Depends on multirow package
% place on the "second" line, and leave space above.
\newcommand{\thead}[1]{
  \multirow{-2}{*}{
    \hspace{-13pt} % RG: Why is there extra space?
    % RG: Get the font and boldface right (reverting to roman)
    \begin{bf}
      \begin{sffamily}
        \begin{tabular}{@{}l@{}}
          #1
        \end{tabular}
      \end{sffamily}
    \end{bf}
  }
}


% Decide whether to include or exclude free-form prose
% Include for internal reviews (e.g. 3rd Year Review); exclude for external CV
%\includecomment{blurb}
\excludecomment{blurb}


% get rid of auto-spacing before and after tables
%\setlength\LTpre{\smallskipamount} % RG: Why doesn't 0pt work?
\setlength\LTpre{0pt}
\setlength\LTpost{0pt}
\newcommand{\tbold}[1]{\textbf{#1}}

\newcommand{\denselist}{\itemsep-3pt\topsep-6pt\partopsep-6pt}
\newcommand{\tailspace}{}

% change sections and subsections: taken from article.cls -- KLB
\makeatletter
\renewcommand\thesection{\@arabic\c@section.\hspace{.19in}}
\renewcommand\thesubsection{\@arabic\c@section.\@arabic\c@subsection.\hspace{.19in}}
\renewcommand\section{\@startsection {section}{1}{\z@}%
  {-3.5ex \@plus -1ex \@minus -.2ex}%
  {2.3ex \@plus.2ex}%
  %                                   {\normalfont\Large\bfseries}}
  {\normalfont\bfseries}}
\renewcommand\subsection{\@startsection{subsection}{2}{\z@}%
  {-3.25ex\@plus -1ex \@minus -.2ex}%
  {1.5ex \@plus .2ex}%
  %                                     {\normalfont\large\bfseries}}
  {\normalfont\itshape}}
\makeatother
% TODO: better would be something like https://tex.stackexchange.com/a/234223/9019

% number subsections alphabetically
\renewcommand\thesubsection{(\alph{subsection})}

% Make "C++" look good
\newcommand{\Cpp}{C\kern-0.05em\texttt{+\kern-0.03em+}}

% Macro
\newcommand{\TODO}[1]{{\textcolor{red}{[TODO: #1]}}}

% Sections that aren't actually there
\newcommand{\nasection}[1]{\addtocounter{section}{1}}
\newcommand{\nasubsection}[1]{\addtocounter{subsection}{1}}


%%%%%%%%%%%%%%%%%%%%%%%%%%%%%%%%%%%%%%%%%%%%%%%%%%%%%%%%%%%%%%%%%%%%%%%%%%%%%%%%
\usepackage[hidelinks]{hyperref}


%% --------------------
%% Bib related
\usepackage[backend=biber,sorting=none,style=numeric,citestyle=numeric,maxbibnames=9999,giveninits=true,defernumbers=true,doi=false,url=false]{biblatex}

% Reduce spacing between references
\setlength\bibitemsep{.2\baselineskip}
%% Show abstract
\DeclareFieldFormat{abstract}{\par\small#1}
\renewbibmacro*{finentry}{\printfield{abstract}\finentry}

%% Bibtex filters
\defbibfilter{journal}{
  keyword=journal and not keyword=noreview and not keyword=submission
}
\defbibfilter{journalnoreview}{
  keyword=journal and keyword=noreview and not keyword=submission
}
\defbibfilter{conf}{
  keyword=conference and not keyword=noreview and not keyword=submission
}
\defbibfilter{confnoreview}{
  keyword=conference and keyword=noreview and not keyword=submission
}
\defbibfilter{workshop}{
  keyword=workshop and not keyword=noreview and not keyword=submission
}
\defbibfilter{workshopnoreview}{
  keyword=workshop and keyword=noreview and not keyword=submission
}
\defbibfilter{noreview}{
  ( keyword=tech-report or keyword=noreview ) and not keyword=submission
}
\defbibfilter{submission}{
  keyword=submission
}

% handle pubstate = {toappear}
% https://tex.stackexchange.com/a/408041/9019
\NewBibliographyString{toappear}
\DefineBibliographyStrings{english}{%
  toappear = {to appear},
}

\renewbibmacro*{in:}{}
%\renewbibmacro*{in:}{%
%  \iffieldundef{pubstate}
%    {}
%    {\printfield{pubstate}%
%     \setunit{\addspace}%
%     \clearfield{pubstate}}%
%  \printtext{%
%    \bibstring{in}\intitlepunct}}

% handle author annotations: equal, mystudent
\newcommand{\bibeqcon}{\ensuremath{^\star}}
\renewcommand*{\mkbibnamegiven}[1]{%
  \ifitemannotation{me}%
    {#1}%{\textbf{#1}}%
    {\ifitemannotation{mystudent}{\textbf{#1}}{#1}}}
\renewcommand*{\mkbibnamefamily}[1]{%
  \ifitemannotation{me}%
    {#1}%{\textbf{#1}}%
    {\ifitemannotation{mystudent}{\textbf{#1}}{#1}}%
  \ifitemannotation{equal}
    {\bibeqcon}
    {}%
}


%%
%% For reverse ordering of references
%% https://tex.stackexchange.com/a/625292/9019

\makeatletter
\newcounter{bibenvcounter}
\newcounter{bibitemcounter}[bibenvcounter]

\newcommand*{\auxkacz@envmax}[2]{%
  \csgdef{blxkacz@envmax@#1}{#2}}

\def\blx@thelabelnumber{%
  \begingroup
  \nottoggle{blx@skiplab}
    {\iftoggle{blx@labelnumber}
       {\ifundef\abx@field@shorthand
         % Need to know if we set any localnumber from .aux file, not
         % just current item. Otherwise, we may just write new
         % localnumbers but not existing ones
         % Also, don't want to regenerate localnumber for any key in a
         % refsection which we've already seen because this has the nasty
         % side-effect of incrementing the localnum counter
         % Also, have to create localnumber again if \nocite{*} as items
         % could have been added to the .bib file without changing the cite
         % information in the .tex
         {\ifboolexpr { (not togl {blx@localnumber}
                         or
                         test {\ifinlist{*}\blx@nocites})
                         and
                         not test {\ifcsdef{blx@defer@\the\c@refsection @\blx@refcontext@context @\abx@field@entrykey}} }
            {\iftoggle{blx@omitnumbers}
               {\let\abx@field@localnumber\@empty}
               {\ifcsundef{blxkacz@envmax@\the\value{bibenvcounter}}
                  {\blx@rerun@latex
                   \csdef{blxkacz@envmax@\the\value{bibenvcounter}}{0}}
                  {\csnumgdef{blx@labelnumber@\the\c@refsection}{%
                   \csuse{blxkacz@envmax@\the\value{bibenvcounter}}%
                      -\value{bibitemcounter}}%
                   \edef\abx@field@localnumber{%
                    \csuse{blx@labelnumber@\the\c@refsection}}%
                  \blx@bbl@addentryfield{\abx@field@entrykey}{\the\c@refsection}%
                  {localnumber}{\blx@refcontext@context}{\abx@field@localnumber}}}}
            {}%
          % Issue re-run message if there was no localnumber in the .aux and
          % \nocite{*} was used because this means that an entry was added to the .bib
          % and was not yet recorded in the .aux. We will have generated a localnumber
          % for such a case above and now we need to re-run to make sure it's in the .aux
          % for the next run
          \ifboolexpr { test {\ifinlist{*}\blx@nocites}
                        and
                        not test {\ifcsdef{blx@localnumber@\the\c@refsection @\abx@field@entrykey}} }
            {\blx@rerun@latex}
            {}%
          % If localnumbers were already in the .aux, add them
          % again from the .aux. This prevents some cycling
          % problems where pagebreaks change after localnumber
          % settles down and then we need another run which then
          % regenerates localnumber requiring another run but then
          % the pagebreaks change back again ... etc.
          \ifcsundef{blx@defer@\the\c@refsection @\blx@refcontext@context @\abx@field@entrykey}
            {\ifundef\abx@field@localnumber
               {}
               {\listxadd\blx@localnumaux{%
                  \string\abx@aux@number%
                     {\the\c@instcount}%
                     {\abx@field@entrykey}%
                     {\the\c@refsection}%
                     {\blx@refcontext@context}%
                     {\abx@field@localnumber}}%
                   % record that we have already generated and output localnum
                   % for this key in this refsection/refcontext
                   \global\cslet{blx@defer@\the\c@refsection @\blx@refcontext@context @\abx@field@entrykey}\@empty}}
           {}}
         {}}%
       {}}
    {}%
  \endgroup}

\defbibenvironment{bibliography}
  {\stepcounter{bibenvcounter}%
   \list
     {\stepcounter{bibitemcounter}%
      \printtext[labelnumberwidth]{%
        \printfield{labelprefix}%
        \printfield{labelnumber}}}
     {\setlength{\labelwidth}{\labelnumberwidth}%
      \setlength{\leftmargin}{\labelwidth}%
      \setlength{\labelsep}{\biblabelsep}%
      \addtolength{\leftmargin}{\labelsep}%
      \setlength{\itemsep}{\bibitemsep}%
      \setlength{\parsep}{\bibparsep}}%
      \renewcommand*{\makelabel}[1]{\hss##1}}
  {\endlist
   \blx@auxwrite\@mainaux{}{%
     \string\auxkacz@envmax
       {\the\value{bibenvcounter}}
       {\the\value{bibitemcounter}}}}
  {\item}
\makeatother



%%
%% Adding Markup to new bib entries (use keyword new)
%%
%% Adapted from here: https://tex.stackexchange.com/questions/411043/biblatex-add-an-image-in-a-bib-entry
\renewbibmacro*{begentry}{%
  \ifkeyword{new}{%
    \new{}%
  }{}%
  \ifkeyword{accepted}{%
    \accepted{}%
  }{}%
  \ifkeyword{changed}{%
    \changed{}%
  }{}%
  \ifkeyword{toappear}{%
    \toappear{}%
  }{}%
}

%%
%% Acceptance rate field
%%
%% Adapted from https://tex.stackexchange.com/questions/111846/biblatex-2-custom-fields-only-one-is-working

\usepackage[many]{tcolorbox} % shaded boxes
\newcommand{\rightsymb}[1]{{% set up
      \parfillskip=0pt % so \par doesnt push \square to left
      \widowpenalty=10000 % so we dont break the page before \square
      \displaywidowpenalty=10000 % ditto
      \finalhyphendemerits=0 % TeXbook exercise 14.32
      %
      % horizontal
      \leavevmode % \nobreak means lines not pages
      \unskip % remove previous space or glue
      \nobreak % don?t break lines
      \hfil % ragged right if we spill over
      \penalty50 % discouragement to do so
      \hskip.2em % ensure some space
      \null % anchor following \hfill
      \hfill % push text to right
      #1% the text to right justify
      %
      % vertical
      \par}} % build paragraph

\colorlet{ARdarkcolor}{black!35!white}
\colorlet{ARlightcolor}{white!96!black}
\colorlet{ARtextcolor}{black!60!white}
\newtcbox{\ARbox}{enhanced,nobeforeafter,tcbox raise base, boxrule=0.4pt, top=-0.4mm, bottom=-0.8mm, right=0mm, left=3.5mm,arc=1pt,boxsep=2pt,before upper={\vphantom{dlg}},   colframe=ARdarkcolor,coltext=ARtextcolor,colback=ARlightcolor,
  overlay={\begin{tcbclipinterior}\fill[ARdarkcolor] (frame.south west) rectangle node[text=white,font=\sffamily\bfseries\tiny,rotate=90]{AR} ([xshift=3.3mm]frame.north west);\end{tcbclipinterior}}}
\newcommand{\ARtext}[1]{%
  \colorlet{saved}{.}\colorlet{ARdarkcolor}{saved!35!white}\colorlet{ARlightcolor}{white!96!saved}\colorlet{ARtextcolor}{saved!60!white}
  \ARbox{\sffamily\footnotesize\textcolor{ARtextcolor}{#1}}}
\newcommand{\AR}[1]{\rightsymb{\ARtext{#1}}}

\DeclareSourcemap{
  \maps[datatype=bibtex]{
    \map{ % turn arinfo into usera for rendering purposes
      \step[fieldsource=arinfo]
      \step[fieldset=usera,origfieldval]
    }
    \map[overwrite]{ % strip eprint (arXiv) info from published papers
      \step[fieldsource=keywords,notmatch=\regexp{submitted|tech-report},final]
      \step[fieldset=eprint,fieldvalue=]
    }
  }
}

% \DeclareFieldFormat{usera}{\sloppy{\ARtext{#1}}
\DeclareFieldFormat{usera}{\ARtext{#1}}

% \AtEveryBibitem{% put at start of entry
% }

\AtEveryBibitem{%
  \ifkeyword{cites1000}{\includegraphics[height=9pt]{figs/1000-cites.pdf}}{}%
  \ifkeyword{cites100}{\includegraphics[height=9pt]{figs/100-cites.pdf}}{}%
  \ifkeyword{cites10}{\includegraphics[height=9pt]{figs/10-cites.pdf}}{}%
  \csappto{blx@bbx@\thefield{entrytype}}{% put at end of entry
    \iffieldundef{usera}{\space}{\space\printfield{usera}}
  }
}

% \newcommand{\remark}[1]{\fontsize{8}{9}\selectfont#1\fontsize{11}{12}\selectfont}
\newcommand{\remark}[1]{{\footnotesize#1}}
\newcommand{\remarkk}[1]{{\footnotesize\par}} % paragraph break so that the paragraph ends with the small font; causes the whole paragraph's baseline to be the small font's baseline
% \newcommand{\remarkk}[1]{\fontsize{8}{9}\selectfont#1\par\fontsize{11}{12}\selectfont} % paragraph break so that the paragraph ends with the small font; causes the whole paragraph's baseline to be the small font's baseline
\newcommand{\tremark}[1]{\vspace{-2em}\par\noindent{\small#1}}


\addbibresource{biblio-cv.bib}

\begin{document}

\nocite{*}

% No page number on the first page
\thispagestyle{empty}

%%
%% Header
%%
\vspace*{-\headsep}\vspace*{-\headheight}
\UBCheader{Curriculum Vitae for Faculty Members}
\vskip 10pt \bigskip\rule{3.416in}{0mm}

\iftoggle{ubcformat}
{
%%
%% Basic info
%%
\noindent
\rule{3.44in}{0mm}%
\parbox[t]{2in}{{\bf Date:} \today}%
\parbox[t]{1in}{Initials: DJS}%
\\[5mm]
\noindent\parbox[t]{.44in}{{\bf 1.}\hfill\mbox{}}%
\parbox[t]{3in}{{\bf SURNAME:} Sutherland}%
\parbox[t]{3.2in}{{\bf FIRST NAME:} Danica}%
\\%
\rule{3.44in}{0mm}%
\parbox[t]{3.45in}{{\bf MIDDLE NAME:} Jane}
\\[5mm]
\parbox[t]{.44in}{{\bf 2.}\hfill\mbox{}}\parbox[t]{6in}{{\bf DEPARTMENT/SCHOOL:} Computer Science}\\[5mm]
\parbox[t]{.44in}{{\bf 3.}\hfill\mbox{}}\parbox[t]{6in}{{\bf FACULTY:} Science} \\[5mm]
\parbox[t]{.44in}{{\bf 4.}\hfill\mbox{}}\parbox[t]{3in}{{\bf PRESENT RANK:}
  Assistant Professor}\parbox[t]{3in}{{\bf SINCE:} December 2020} \\[5mm]
\setcounter{section}{4}

%%
%% 5. POST-SECONDARY EDUCATION
%%
\vspace{-30pt}
\ulsection{POST-SECONDARY EDUCATION}
% \subsection{Degrees}
% \noindent
\noindent\vspace{-2em}
\begin{table}[h!]
  % \bf
  % \begin{tabular}{||p{0.4\textwidth}||>{\centering\let\newline                                  \\\arraybackslash\hspace{0pt}}p{0.1\textwidth}||>{\centering\let\newline                                  \\\arraybackslash\hspace{0pt}}p{0.2\textwidth}||>{\centering\let\newline                                  \\\arraybackslash\hspace{0pt}}p{0.2\textwidth}||}
  % \hhline{#=#=#=#=#}
  % University or Institution                  & Degree & Subject Area     & Dates \\
  % \hhline{#=#=#=#=#}
  \begin{tabular}{ p{0.3\textwidth} p{0.1\textwidth} p{0.28\textwidth} p{0.20\textwidth} }
    \toprule
    \multicolumn{1}{c}{\bf University or Institution} & \multicolumn{1}{c}{\bf Degree} & \multicolumn{1}{c}{\bf Subject Area}            & \multicolumn{1}{c}{\bf Dates} \\
    \midrule
    Carnegie Mellon University & PhD + MS & Computer Science & 2011/08 -- 2016/09 \\
    Swarthmore College         & BA       & Computer Science & 2007/08 -- 2011/05 \\
    \bottomrule
  \end{tabular}
\end{table}

% \hspace{1.8em}\dunderline{1pt}{\textbf{Special Professional Qualifications}}

%%
%% 6. EMPLOYMENT RECORD
%%
\ulsection{EMPLOYMENT RECORD}
\subsection{Prior to coming to UBC}
\noindent\vspace{-2em}
\begin{table}[h!]
  % \bf
  % \begin{tabular}{||p{0.4\textwidth}||>{\centering\let\newline                                  \\\arraybackslash\hspace{0pt}}p{0.1\textwidth}||>{\centering\let\newline                                  \\\arraybackslash\hspace{0pt}}p{0.2\textwidth}||>{\centering\let\newline                                  \\\arraybackslash\hspace{0pt}}p{0.2\textwidth}||}
  % \hhline{#=#=#=#=#}
  % University or Institution                  & Degree & Subject Area     & Dates \\
  % \hhline{#=#=#=#=#}
  \begin{tabular}{ p{0.4\textwidth} p{0.3\textwidth} p{0.20\textwidth} }
    \toprule
    \multicolumn{1}{c}{\bf University, Company or Organization} & \multicolumn{1}{c}{\bf Rank or Title} & \multicolumn{1}{c}{\bf Dates} \\
    \midrule
    Toyota Technological Institute at Chicago  & Research Assistant Professor  & 2019/12 -- 2020/08 \\
    University College London & Research Associate   & 2016/10 -- 2019/08  \\
    \bottomrule
  \end{tabular}
\end{table}


\vspace{-1em}
\subsection{At UBC}
\noindent\vspace{-2em}
\begin{table}[h!]
  % \bf
  \begin{tabular}{ p{0.725\textwidth} p{0.20\textwidth} }
    \toprule
    \multicolumn{1}{c}{\bf Rank or Title} & \multicolumn{1}{c}{\bf Dates} \\
    \midrule
    Canada CIFAR AI Chair, Alberta Machine Intelligence Institute & 2021/06 -- present \\
    Assistant Professor               & 2020/12 -- present        \\
    \bottomrule
  \end{tabular}
\end{table}

% Not yet!
\nasubsection{%
  \vspace{-1em}
  \subsection{Date of granting of tenure at U.B.C.: \textmd{N/A}}
}%

{%\nasection{%
  %
  % 7. LEAVES OF ABSENCE
  %
  \ulsection{LEAVES OF ABSENCE}
  \noindent\vspace{-2em}
  \begin{table}[h!]
    \begin{tabular}{ p{0.45\textwidth} p{0.25\textwidth} p{0.20\textwidth} }
      \toprule
      \multicolumn{1}{c}{\parbox[c][][c]{0.45\textwidth}{\centering\bf University, Company or Organization \\ at which Leave was taken}} &
      \multicolumn{1}{c}{\bf Type of Leave}
      & \multicolumn{1}{c}{\bf Dates} \\
      \midrule
      University of British Columbia &  Medical & 2022/06 -- 2022/08 \\
      \bottomrule
    \end{tabular}
  \end{table}

  \clearpage
}%

}{
    \begin{center}
    \textbf{\Large Danica J.\ Sutherland}
    \\
    \footnotesize{\textit{Version as of \today.}}
    \end{center}

\ulsection{ACADEMIC POSITIONS AND EDUCATION}
\begin{table}[h!]
  \begin{tabular}{ p{0.36\textwidth} p{0.40\textwidth} p{0.13\textwidth} }
    \toprule
    %\multicolumn{1}{c}{\bf Rank or Title} & \multicolumn{1}{c}{\bf Dates} \\
    %\midrule
    Assistant Professor, Computer Science & University of British Columbia  & 2020 --         \\
    Canada CIFAR AI Chair & Alberta Machine Intelligence Institute & 2021 --  \\
    \midrule
    Research Assistant Professor & Toyota Technological Institute at Chicago & 2019 -- 2020 \\
    Research Associate & University College London & 2016 -- 2019 \\
    Ph.D., Computer Science & Carnegie Mellon University & 2011 -- 2016 \\
    B.A., Computer Science & Swarthmore College & 2007 -- 2011 \\
    \bottomrule
  \end{tabular}
\end{table}
}

%%
%% 8. TEACHING
%%
\ulsection{TEACHING} \label{sec:teaching}

\iftoggle{ubcformat}
{
\subsection{Areas of special interest and accomplishments}
% UBC rules: 150 words max

  \paragraph{Design of CPSC 532D: Modern Statistical Learning Theory.}
  \continuing{}
  This graduate course
  %(in a mostly lecture-oriented format, due to the quantity of necessary content) focuses on
  giving students the broad background in learning theory
  required for mathematical understanding of machine learning methods.
  The course integrates material from several textbooks and other courses
  with multiple lectures of material from research papers published in the last five years.
  %I combined materials from several textbooks and notes sets with significant extra material from current research papers to not only cover traditional topics developed in the 70s and 80s, but also bring them up to speed on current developments, with multiple lectures devoted almost exclusively to papers published since 2018.
  %\changed{In the second offering this year, I significantly further developed the course, including a format change (from slides to board with lecture notes) and changes to the presentation and included material.}
  My lecture notes, much-appreciated by the students,
  cover these varied topics in a way that strives to be accessible while also correcting mistakes in the literature.

  \paragraph{Update of CPSC 440/550: Advanced Machine Learning}
  \changed{I am the first person other than Mark Schmidt to teach the current form of this course.
  I have continued to modernize the topics covered in the course and the assignments
  (switching to further integrated use of modern standard tools),
  as well as pedagogical changes such as introducing biweekly quizzes with the Computer-Based Testing Facility (I believe the first AI course to do so).}

%  \paragraph{CPSC 340: Machine Learning and Data Mining.}
%    I also \new{taught CPSC 340 for the first time this year}
%    (running everything except lecturing in conjunction with Mike Gelbart's section).
%    This is a mature course whose materials I mostly only tweaked,
%    although I modernized some aspects of the assignments in particular,
%    and planned for some more substantial changes in my next offering.
}{
}

\subsection{Courses Taught at UBC}

\iftoggle{ubcformat}{}{
    ~\vspace*{-3ex}

    CPSC 532D: Modern Statistical Learning Theory
    is a graduate course giving students the broad background in learning theory
    required for mathematical understanding of machine learning theory.

    CPSC 440/540: Advanced Machine Learning
    is a second course in machine learning,
    covering a variety of topics with a probabilistic bent,
    intended for advanced undergraduates or beginning graduate students.

    CPSC 340: Machine Learning and Data Science
    is an introductory course in machine learning.
}

\begin{coursetable}
  \newtable 23-24 W2  & CPSC 440+550 & 3 & 105+25        & 3/week  & 1h/week office hours \\
  \newtable 23-24 W1  & CPSC 532D    & 3 & 15            & 3/week  & 2h/week office hours \\
            22-23 W2  & CPSC 440+540 & 3 & 70+30         & 3/week  & 1h/week office hours \\
            22-23 W1  & CPSC 532D    & 3 & 23            & 3/week  & 1h/week office hours \\
            21-22 W2  & CPSC 532S    & 3 & 23            & 3/week  & 2h/week office hours \\
            21-22 W1  & CPSC 340     & 3 & 89\textit{\small +139} & 3/week  & 1h/week office hours \\
\end{coursetable}
\tremark{%
CPSC 440 and 550 (previously 540) are undergraduate and graduate cross-listings of one course; sizes show the enrolment in each listing.\\
CPSC 532D was formerly known as 532S.\\
CPSC 340 in 2021-22 W1 was run jointly with another section taught by Mike Gelbart with an enrolment of 139, for a total of 228 students across the two sections.\\
}

% Course number details:
% \begin{itemize}[nosep]
%   \item CPSC 340: Machine Learning and Data Mining
%   \item CSC 532S: Topics in AI: Modern Statistical Learning Theory
% \end{itemize}

\subsection{Graduate Research Supervision}
% \subsection{Student Supervision}
% All activities below the dashed line are supervisory activities that ended prior to joining UBC.

% \paragraph{Postdoc Supervision}
% \begin{suptable}
%   \newtable Fabrizio Pedersoli & Postdoc UBC  & 2021/01 & ongoing & Supervisor \\
%   \hdashline
%   Baptiste Angles    & Postdoc UVic & 2019/04 & 2019/08 & Supervisor \\
% \end{suptable}
% \vspace{-2em}
%
% \noindent
% {\small Baptiste now works at Apple.}

%\paragraph{PhD Supervision}
\paragraph{Supervision}\footnotetext[1000]{Switched to my supervision at this point.}% XXX idk why this is necessary, blegh
\begin{suptable}
  \newtable        Bingshan Hu   & Postdoc, UBC CS & 2023/11 & ongoing & Primary supervisor \\
  \newtable        Nathaniel Xu  & PhD, UBC CS & 2023/09 & ongoing & Supervisor \\
  \newtable        Zheng He      & PhD-track, UBC CS & 2023/09 & ongoing & Supervisor \\
  \continuingtable Hamed Shirzad & PhD, UBC CS & 2021/09 & ongoing & Supervisor \\
  \continuingtable Wonho Bae     & PhD, UBC CS & 2020/09 & ongoing & Supervisor \\
  \continuingtable Yi Joshua Ren & PhD, UBC CS & 2020/09 & ongoing & Supervisor \\
%\end{suptable}
%
%\paragraph{MSc Thesis Supervision}
%\begin{suptable}
  \changedtable Mohamad Amin Mohamadi & MSc, UBC CS & 2021/05\switched{} & 2023/04 & Supervisor  \\
  \newtable                           & Researcher  & 2023/05 & 2024/01 & Supervisor  \\
  \changedtable Namrata Deka          & MSc, UBC CS & 2020/09 & 2023/04 & Supervisor    \\
  \newtable                           & Researcher  & 2023/05 & 2023/08 & Supervisor  \\
  \changedtable Milad Jalali          & MSc, UBC CS & 2020/09 & 2023/08 & Supervisor    \\
\end{suptable}
\tremark{%
  Bingshan is co-supervised with Mathias Lécuyer and Trevor Campbell; I provide the majority of her funding.
  Mohamad Amin and Namrata, after graduating,
  continued for some time in non-student research positions,
  before beginning their PhDs at
  TTI-Chicago and Carnegie Mellon, respectively.
  Milad is currently seeking an industry job.
}
%  \hdashline
%  \continuingtable \textit{Lijia Zhou} & \textit{PhD, UChicago Stat} & \textit{2020/03} & \textit{ongoing} & \textit{Informal co-supervisor}
%  \\
%  % \hdashline
%  \textit{Akilesh Tangella}  & \textit{MSc, TTI-Chicago} & \textit{2020/03} & \textit{2020/12} & \textit{Informal co-supervisor}
% \tremark{%
%   UChicago Statistics does not allow me a formal supervisory role or even a committee position, but I have weekly advising meetings with Lijia.
%   Akilesh moved on to a Research Engineer position at Webb.
% }

\iftoggle{ubcformat}{\clearpage}{} % XXX TEMP
\paragraph{Graduate Supervisory Committee Memberships}
\begin{supcommtable}
  \continuingtable Victor Portella        & PhD, UBC CS       & 2023/01 & ongoing & PhD committee member \\
  \newtable        Chen Fan               & PhD, UBC CS       & 2023/05 & 2023/07 & RPE committee member \\
  \newtable        Saiyue Lyu             & PhD, UBC CS       & 2023/05 & 2023/07 & RPE committee member \\
  \newtable        Qiaoyue Tang           & PhD, UBC CS       & 2023/04 & 2023/06 & RPE committee member \\
                   Amrutha Ramesh         & PhD, UBC CS       & 2022/09 & 2022/11 & RPE committee member \\
                   Melika Farahani        & MSc, UBC CS       & 2022/06 & 2022/08 & MSc committee member \\
\end{supcommtable}
\tremark{%
  The RPE, ``research proficiency exam,'' is a step in the degree for PhD students at UBC CS roughly corresponding to a qualifying exam.
}


\subsection{Graduate Program Supervision}
\begin{suptable}
  \newtable Anubhav Garg  & MSc (essay), UBC CS & 2022/10\switched{} & ongoing & Supervisor \\
  \newtable Arsh Jhah     & MSc (essay), UBC CS & 2022/06\switched{} & 2023/10 & Supervisor \\
\end{suptable} \label{sec:other-sup}
\tremark{%
  Arsh and Anubhav both came to me as essay-option students, switching from other advisors, rather than switching from a research program under my supervision. {\color{white} \switched}
}


\subsection{Continuing Education Activities} \label{sub:continuing-ed}

\begin{itemize}
    {%- for talk in talks if talk.teaching -%}
{%- set venue = venues[talk.venue] -%}

{%- if not loop.previtem or (talk.date.year != loop.previtem.date.year) %}
\item[\textbf{ {{- talk.date.year -}} }]
{% else -%}
\\
{% endif -%}
  {% if talk|in_last_year(month=5) -%} \new{ {%- endif -%}
    \textbf{ {{- talk.talkmeta or "Instructor" -}} },
      {{ venue.name -}}
      {{- venue.short | maybe_wrap(" (", ")") -}}
      {%- if venue.conference -%}
        {%- set conf = venues[venue.conference] -%}
        {{- "" }} (at {{ conf.long_article | maybe_wrap("", " ") -}} {{- conf.name -}} {{- conf.short | maybe_wrap(", ", "") }})
      {%- endif -%}
      .
      {%- if talk.funded -%} \paid{} {%- endif -%}
      {%- if talk.competitive -%} \competitive{} {%- endif -%}
      {%- if talk.virtual -%} \virtual{} {%- endif %}
      \remark{\emph{ {{- talk.title -}} }.
              {{- talk.note | maybe_wrap(" ", "") -}}
              {{- talk.cv_note | maybe_wrap(" ", "") -}} }

  
  {%- if talk|in_last_year(month=5) -%} } {%- endif -%}
{%- endfor -%}

{# 
{% for paper in papers if paper.my_talk -%}
{%- set venue = venues[paper.venue] %}

\item[ {{- paper.year -}} ]
  {% if paper|in_last_year(month=5) %}\new{ {%- endif -%}
  {{- venue.name -}}
  {{- venue.short | maybe_wrap(" (", ")") -}}
  {%- if venue.conference -%}
    {%- set conf = venues[venue.conference] -%}
    {{- "" }} (workshop at {{ conf.long_article | maybe_wrap("", " ") -}} {{- conf.name -}} {{- conf.short | maybe_wrap(", ", "") }})
  {%- endif -%}
  .
  \remark{\emph{  {{- paper.title | maybe_punc(".") -}}  }}
  {%- if paper|in_last_year(month=5) -%} } {%- endif -%}
{% endfor %}
#}

\end{itemize}


\nasubsection{%
  \subsection{Visiting Lecturer (indicate university/organization and dates)}

  Not applicable.
}%


\nasubsection{%
  \subsection{Other}

  Not applicable.
}%

%%
%% 9. SCHOLARLY AND PROFESSIONAL ACTIVITIES
%%
\ulsection{SCHOLARLY AND PROFESSIONAL ACTIVITIES}
\iftoggle{ubcformat}{
\subsection{Areas of special interest and accomplishments} % UBC rules: 150 words max

% My lab's work broadly spans representation learning, statistical testing, learning theory, and understanding probability distributions. A few recent focus areas:

%Work accepted in the last year focuses on two areas:
%
\paragraph{Understanding modern machine learning.}
\continuing{
Many aspects of my lab's work are based on improving understanding of how modern deep learning methods work.
One particular trend is in new uses for neural tangent kernels,
including improved tools for their empirical study \parencite{mohamadi:pseudo-ntk},
applications to theoretical analysis \parencite{ren:finetuning,ren:learning-path},
and practical uses \parencite{mohamadi:active-ntk}.
Work with related tools has also contributed to understanding the grokking phenomenon (a hotly-studied problem in learning theory) \parencite{mohamadi:grokking},
and to deepen foundational understanding of learning theory in high-dimensional regimes
\parencite{zhou:moreau,zhou:optimistic-rates,koehler:gaussian-interpolators}.
}

\paragraph{Improved techniques based on theoretical insights.}
Our theoretical insights have also led to improved models for
learning on graph data (such as molecular graphs and citation networks) \parencite{shirzad:exphormer,shirzad:sparsifying-exphormer},
private generative models \parencite{harder:dp-mepf,yang:dp-ntk},
weakly-supervised semantic segmentation \parencite{bae:prediction-filtering},
and the closely related problems of improving out-of-distribution generalization and fairness \parencite{ren:sem-il,pogodin:circe,deka:mmd-bfair}.

%
%\paragraph{Kernels for improved deep learning.}
%\continuing{
%Integrating kernel methods into deep learning models
%has provided significant benefits across a range of machine learning contexts,
%and helps various attractive and theoretically grounded kernel models work well in practice.
%Recent progress includes
%methods to enforce conditional invariance constraints
%(highly applicable to fairness and out-of-distribution generalization settings),
%in both continuous \parencite{pogodin:circe} and more traditional discrete \parencite{deka:mmd-bfair} settings.
%% as well as improved methods for evaluating graph generative models \parencite{shirzad:contrastive-graph-eval}.
%}
}{}


\subsection{Research or equivalent grants}

% \paragraph{Since joining UBC}
% \vspace{-0.5em}
% \begin{itemize}
%   \setlength\itemsep{-0.5em}
%   \item Total (co/)awarded: \new{\$461,442}
%   \item Research grants: \new{\$311,190}
%   \item Compute cluster cost: \new{\$150,252}
% \end{itemize}

\begin{granttable}
\iftoggle{ubcformat}{
  \newtable NFRF
  & Automated Scanning Probe Microscopy for Investigating Surface Chemistry (Application under review)
  & \$125,000 & 2024--26 &
  S.~Burke; \textit{J.~Hein}; \textit{A.~Schiffrin}; \textit{D.~J.~Sutherland}
  \\ \midrule
  \newtable NSERC, CSE
  & URGENCY: foUndations foR scalinG ai sEcurity aNd privaCY (Application under review)
  & $\!\!\!\!\!$\$1,400,000 & 2024-28 &
  F.~Kirschbaum; \textit{U.~Aivodji}; \textit{N.~Asokan}; \textit{W.~Chen}; \textit{C.~Clarke}; \textit{S.~Gambs}; \textit{A.~Gurfinkel}; \textit{X.~He}; \textit{G.~C.~Kamath}; \textit{M.~Lécuyer}; \textit{M.~Smucker}; \textit{V.~Shwartz}; \textit{D.~J.~Sutherland}; \textit{T.~Vidal}; \textit{Y.~Yu}
  \\ \midrule\midrule
}{}
  \newtable
  DRAC, Resources for Research Groups
  & Empirically Understanding Deep Learning with Theory\competitive{} %&C
  & \$24,449\compequiv & 2024--25
  & D.~J.~Sutherland
  \\ \midrule
  \continuingtable
  NSERC, CREATE
  & Advanced Machine Learning Training Network\competitive %&C
  & $\!\!\!\!\!$\$1,075,650\footnote{Includes contributions from various UBC departments and partners; NSERC amount is on average \$275,000/year.} & 2023-29
  & F.~Wood; \textit{{A.~Bouchard-Côté;} {B.~Bloem-Reddy;} {T.~Campbell;} {K.~Leyton-Brown;} {M.~van de Panne;} {D.~Poole;} {M.~Schmidt;} {L.~Sigal;} {D.~J.~Sutherland}; {K.~M.~Yi}}
  \\ \midrule
  \continuingtable
  NSERC + FRQNT, NOVA
  & Towards Improved Understanding and Efficient Utilization of Depthwise Computation in Modern Neural Networks\competitive{} %& C
  & \$75,000 & 2023--26
  & G.~Wolf; \textit{E.~Belilovsky; D.~J.~Sutherland; J.~Tian}
  \\ \midrule
  %\newtable
  UBC Data Science Institute, Postdoctoral Matching Fund
  & Anytime-Valid PAC-Bayes for Industrial Applications\competitive{} %& C
  & \$40,000 & 2023--24
  &  T.~Campbell; \mbox{D.~J.~Sutherland}
  \\ \midrule
  %\newtable
  DRAC, Resources for Research Groups
  & Empirically Understanding Deep Learning with Theory\competitive{} %&C
  & \$28,737\compequiv & 2023--24
  & D.~J.~Sutherland
  \\ \midrule
  \continuingtable
  CIFAR  & Canada CIFAR AI Chair, Amii\competitive{} % & C
  & \$78,333 & 2021--26 & D.~J.~Sutherland
  \\ \midrule
  \continuingtable
  NSERC Discovery & Learning Better Representations with Kernels\competitive{} % & C
  & \$29,000 & 2021--26 & D.~J.~Sutherland
  \\ \midrule
  NSERC Discovery & Launch Supplement % & NC
  & \$12,500 & 2021 & D.~J.~Sutherland
  \\ \midrule
  UBC  & Startup Fund %& NC 
  & \$200,000 & 2020  & D.~J.~Sutherland \\
\end{granttable}
%\tremark{Under the header ``Comp'', C refers to competitive grants and NC to non-competitive.}

\nasubsection{%
  \subsection{Research or equivalent contracts}
  Under the header ``Comp'', C refers to competitive grants and NC to non-competitive.
  \begin{granttable}
    N/A &  &  &  &   & \\
  \end{granttable}
}%
\footnotetext[1003]{Estimated equivalent cost for computational resources.} % XXX

\subsection{Invited Presentations} % (Identify whether International/National/Local)}

Invited tutorials and summer school lectures are listed instead in section \ref{sub:continuing-ed}.
% TODO: footnote symbols getting kinda silly; figure out 
{\color{white} \competitive\international\national\local\virtual\paid}

\begin{itemize}[leftmargin=5em]
{% for talk in talks if not talk.teaching -%}
{%- set venue = venues[talk.venue] %}

\item[ {{- "{:%b %Y}".format(talk.date) -}} ]
  {% if talk|in_last_year(month=5) %}\new{ {%- endif -%}
  {{- venue.name -}}
  {{- venue.short | maybe_wrap(" (", ")") -}}
  {%- if venue.conference -%}
    {%- set conf = venues[venue.conference] -%}
    {{- "" }} (workshop at {{ conf.long_article | maybe_wrap("", " ") -}} {{- conf.name -}} {{- conf.short | maybe_wrap(", ", "") }})
  {%- endif -%}
  .
  {#- talk.type | maybe_wrap("$^\\text{", "}$ ") -#}
  {{- talk.type | maybe_wrap("\\", "{}") -}}
  {%- if talk.funded %}\paid{}{% endif -%}
  {%- if talk.virtual %}\virtual{}{% endif -%}
  {{- " " -}}\remark{
    {{- "" -}}
    \emph{ {{- talk.title | maybe_punc(".") -}} }
    {{- talk.note | maybe_wrap(" ", "") -}}
    {{- talk.cv_note | maybe_wrap(" ", "") -}}
    {#- talk.type | capitalize | maybe_punc(".") -#}
  }
  {%- if talk|in_last_year(month=5) -%} } {%- endif -%}
{% endfor %}

\end{itemize}


\subsection{Other Presentations}

\begin{itemize}[leftmargin=5em]
{% for paper in papers if paper.my_talk -%}
{%- set venue = venues[paper.venue] %}

\item[ {{- paper.year -}} ]
  {% if paper|in_last_year(month=5) %}\new{ {%- endif -%}
  {{- venue.name -}}
  {{- venue.short | maybe_wrap(" (", ")") -}}
  {%- if venue.conference -%}
    {%- set conf = venues[venue.conference] -%}
    {{- "" }} (workshop at {{ conf.long_article | maybe_wrap("", " ") -}} {{- conf.name -}} {{- conf.short | maybe_wrap(", ", "") }})
  {%- endif -%}
  .
  \remark{\emph{  {{- paper.title | maybe_punc(".") -}}  }}
  {%- if paper|in_last_year(month=5) -%} } {%- endif -%}
{% endfor %}

\end{itemize}


\nasubsection{%
  \subsection{Other}

  Not applicable.
}%

\subsection{Conference Participation} % (Organizer, Keynote Speaker, etc.)
\label{sec:conference-participation}

Reviewing-based roles are in section \ref{sec:reviewer}; tutorials are in section \ref{sub:continuing-ed}.

\begin{itemize}[leftmargin=5em]
  \item[2020-24] \changed{Area Chair, Conference on Artificial Intelligence and Statistics (AISTATS).}
  \item[2020-22] \changed{Area Chair, Conference on Neural Information Processing Systems (NeurIPS).}
  \item[2022] {Panel Moderator, \emph{Faculty and Queerness}, Queer in AI NeurIPS workshop.}
  \item[2022] {Area Chair, International Conference on Machine Learning (ICML).}
  \item[2022] {D\&I event organizer (Queer in AI), AAAI Conference on Artificial Intelligence (AAAI).}
  \item[2021-22] {Senior Program Committee, AAAI Conference on Artificial Intelligence (AAAI).}
\end{itemize}



%%
%% 10. SERVICE TO THE UNIVERSITY
%%
% \ulsection{SERVICE TO THE UNIVERSITY AND THE COMMUNITY}

% Please include a clear and concise statement of your service (including any service undertaken to advance the inclusion of all those who have been historically excluded based on gender, race, religion, sexuality, age, disability, or economic circumstance) and the criteria you deem to be appropriate in assessing your contributions. (Statements exceeding 150 words will not be considered).

% \textbf{TODO}


% TODO: redefine section* properly
% \section*{\hspace{.52in}\dunderline{1pt}{SERVICE TO THE UNIVERSITY}}
\ulsection{SERVICE TO THE UNIVERSITY}


\subsection{Areas of special interest and accomplishments}

% \paragraph{Beyond deep learning.}
% \new{%
% I am involved in efforts led by Frank Wood to form an industrial consortium.
% }


% \subsection{Areas of service undertaken to advance the inclusion of all those who have been historically excluded based on gender, race, religion, sexuality, age, disability, or economic circumstance}

% \paragraph{Women in Science mentorship}
% \new{In 2023, I began participating in the UBC Women in Science mentorship program.}
% I currently meet with two first-year undergraduate women to give them advice and mentorship on their careers in science.

\paragraph{Informal mentorship for trans undergrads}
\continuing{Since beginning at UBC, I've occasionally met with several trans undergrads in Computer Science and related fields for mentorship, career advice, and community.}

\paragraph{Preferred names.}
{%
Many UBC IT systems ignore preferred names,
a significant problem for e.g.\ trans people who have not or cannot change their legal name in BC.
I made a list of over a dozen in 2021-22,
some of which I've worked directly towards getting fixed,
some of which I've just handed off to the Equity Office or others.
Since I changed my legal name in 2022, it's been harder for me to find more instances.
}

\paragraph{GRAC process improvements.}
In 2020-21, I helped streamline processes (including writing code) for coordinating internal review of AI-area graduate applicants.


\subsection{Memberships on committees} % including offices held and dates

\begin{itemize}[leftmargin=5em]
  \item[2023--24] \new{Faculty recruiting committee: Artificial Intelligence Methods for Scientific Impact.}
  \item[2023] \new{Ad-hoc faculty recruiting committee (single candidate).}
  \item[2023] AIM-SI Steering Committee (temporary position as a leave replacement).
  \item[2022--23] Graduate Recruiting and Admissions Committee, Computer Science Department.
  \item[2021--22] Faculty recruiting committee: CRC Tier 2 in Quantum Computing, Computer Science Department.
  \item[2021] Graduate Recruiting and Admissions Committee, Computer Science Department.
\end{itemize}

%\nasubsection{%
\subsection{Other service} % including dates

\begin{itemize}[leftmargin=5em]
  \item[2023--] \new{Recruitment Chair, Advanced Machine Learning Training Network.}
  \item[2023] {UBC Women in Science mentorship program}, occasionally meeting first-year undergraduates.
  \item[2024] Chair of thesis proposal exam: \new{William Harvey}.
  \item[2023] Chair of thesis proposal exam: Christian Weilbach, Shih-Yang Su.
  \item[] Chair of thesis defense: Andreas Munk.
  \item[2022] Chair of thesis proposal exam: Vaden Masrani.
\end{itemize}
%}%


%%
%% 11. SERVICE TO THE COMMUNITY
%%
\ulsection{SERVICE TO THE COMMUNITY}

% \subsection{Service undertaken to advance the inclusion of all those who have been historically excluded based on gender, race, religion, sexuality, age, disability, or economic circumstance}

\nasubsection{
\subsection{Memberships on scholarly societies %, including offices held and dates%
           }
\begin{itemize}[leftmargin=5em]
  % \item[2018/10--Now] Adjunct Member, Astronomy Research Centre (University of Victoria)
\end{itemize}
}

%\nasubsection{%
\subsection{Memberships in other societies} %, including offices held and dates}
\begin{itemize}[leftmargin=5em]
\item[2022--] \continuing{Queer in AI, core organizer.} \remark{%
  (Also various capacities from 2018 to 2021.)
  My work with this international nonprofit organization has included
  trans-inclusive advocacy to conferences, particularly making enormous improvements to the way in which NeurIPS, ICML, ICLR, AISTATS, and MLSys handle attendees with professional names different from their legal names;
  working to improve post-publication name change policies at several venues;
  co-organizing social, mentoring, and workshop events at conferences;
  serving on a hiring committee for paid administrator positions;
  initiating and maintaining an informal program of feedback and mentoring for faculty candidates;
  and more.
}
% \paragraph{Queer in AI.}
% I am involved with a nonprofit organization called Queer in AI,
% \continuing{serving as a core organizer since January 2022.}
% My work through the organization has included
% trans-inclusive advocacy to conferences
% (including enormous improvements to handling of preferred names at NeurIPS/ICML/ICLR/AISTATS/MLSys and \continuing{improved post-publication name change policies at several venues});
% {co-organizing social, mentoring, and workshop events at conferences};
% reviewing administrator applications as part of a hiring committee;
% \continuing{initiating an informal program of feedback and mentoring for faculty candidates};
% and more.
\item[2021--] \continuing{Name Change Policy Working Group, member.} \remark{This loosely-organized group (no formal positions) advocates with scientific publishers for post-publication name change policies. I gave two invited talks on these issues in 2022-23.}
\end{itemize}
% Not applicable.
%}%

\nasubsection{%
  \subsection{Memberships on scholarly committees} %, including offices held and  dates}

  Not applicable.
}%

%\nasubsection{%
\subsection{Memberships on other committees} %, including offices held and dates}
\begin{itemize}[leftmargin=5em]
\item[2022-23] {CIFAR Pan-Canadian AI Strategy's National Program Committee, member.}
\end{itemize}
% Not applicable.
%}%

\nasubsection{%
  \subsection{Editorships (list journal and dates)}

  Not applicable.
}%

\subsection{Reviewer} \label{sec:reviewer} % (journal, agency, etc., including dates)

Area chair / meta-reviewer roles are listed under section \ref{sec:conference-participation}.

\begin{itemize}[leftmargin=5em]
  \item[2024] Conference reviewing: ICML.
  \item[] Grant reviewing: CIFAR AI Catalyst grant.
  \item[2023] Journal reviewing: {JRSS-A}.
  \item[    ] Conference reviewing: {ICLR, NeurIPS}.
  \item[] Grant reviewing: NSERC Discovery Grant.
  \item[2022] Journal reviewing: {JMLR}.
  \item[    ] Conference reviewing: {UAI}, {ICML}.
  \item[    ] Grant reviewing: NSERC Discovery Grant.
  \item[2021] Journal reviewing: {JMLR}, {Bernoulli}.
  \item[] Conference reviewing: ICML.
  \item[2014--20] Conference reviewing: NeurIPS, ICML, ICLR, AISTATS, AAAI, COLT, SoCG, IJCAI, ECML-PKDD.
  \item[] Journal reviewing: JMLR, IEEE TSP, IEEE T-PAMI, Bernoulli, MLJ, Comptes rendus.
\end{itemize}


\iftoggle{ubcformat}{\clearpage}{} % XXX TEMP
\subsection{External examiner} % (indicate university and dates)}

\begin{extcomtable}
  2020 & Iryna Korshunova & Ghent University & Doctor of Computer Science Engineering \\
  2020 & Tong Che & University of Cambridge & Master of Philospohy \\
\end{extcomtable}

\nasubsection{
\subsection{Consultant} % (indicate organization and dates)}
}

\nasubsection{
\subsection{Other service to the community}
Not applicable.

% \paragraph{Queer in AI.}
% I am involved with a nonprofit organization called Queer in AI,
% \continuing{serving as a core organizer since January 2022.}
% My work through the organization has included
% trans-inclusive advocacy to conferences
% (including enormous improvements to handling of preferred names at NeurIPS/ICML/ICLR/AISTATS/MLSys and \continuing{improved post-publication name change policies at several venues});
% {co-organizing social, mentoring, and workshop events at conferences};
% reviewing administrator applications as part of a hiring committee;
% \continuing{initiating an informal program of feedback and mentoring for faculty candidates};
% and more.

% \paragraph{Name Change Policy Working Group}
% \continuing{I am also a member of an informal group that advocates for post-publication name change policies,}
% and gave two invited talks on these issues in the last year.

}

%%
%% 12. AWARDS AND DISTINCTIONS
%%
\ulsection{AWARDS AND DISTINCTIONS}

%\nasubsection{%
\subsection{Awards for Teaching} % (indicate name of award, awarding organizations and date)

\begin{itemize}[leftmargin=5em]
    \item[] Letters from the Dean of the Faculty of Science for ``some of the highest student evaluations in the Faculty of Science'' for: CPSC 340 in 2021W1, CPSC 532S in 2021W2.
\end{itemize}
%}%

\subsection{Awards for Scholarship} % (indicate name of award, awarding organizations and date)

\begin{itemize}[leftmargin=5em]
  \item[2021] {Canada CIFAR AI Chair, Amii. \remark{From CIFAR: ``The goal of the Chairs Program is to recruit and retain in Canada some of the world's leading researchers in AI and provide them with long-term, dedicated research funding to support their research programs and help them train the next generation of AI leaders.'' Accompanied by funding, listed under grants.}} %https://www.cifar.ca/ai/pan-canadian-artificial-intelligence-strategy

  \item[2014--16] Sandia Campus Executive Program. \remark{A PhD fellowship, renewed in 2015.}
  \item[2013] National Science Foundation Graduate Research Fellowship: Honorable mention.
  \item[2011] Undergraduate awards: \remark{Ivy Award, one of two awards to the graduating class for being ``outstanding in leadership, scholarship, and contributions to the college community.'' Elected Phi Beta Kappa.}
\end{itemize}

Additionally, the following papers received best paper awards,
or competitive presentations selected by area and program chairs at large conferences
which might be thought of as comparable to best paper awards at a smaller venue:
% Additionally, the following papers were selected for competitive oral or spotlight presentations by area and program chairs,
% and thus might be thought of as comparable to an outstanding paper award at a small venue:
\begin{itemize}[leftmargin=5em]
  % TODO: make this a table, to right-align the citation keys
  {%- for paper in papers if paper["accept-kind"]
    and (venues[paper.venue].type == "conference" or paper["accept-kind"] == "best paper") -%}
{%- set venue = venues[paper.venue] %}

\item[
  {%- if not loop.previtem or (paper.year != loop.previtem.year) -%}
    {{- paper.year -}}
  {%- endif -%}
]
  {% if paper|in_last_year(month=5) -%} \new{ {%- endif -%}
  \cite{ {{- paper | bibtex_key(coauthors) -}} }:
%
  {{ venue.short or venue.name }}
  {{ paper["accept-kind"] | bibtex_escape }}:
  {{ paper | ar_info(venue, verbose=True, first_only=True) | bibtex_escape }}.
  {% if paper|in_last_year(month=5) -%} } {%- endif -%}

{%- endfor -%}

\end{itemize}

\subsection{Awards for Service} % (indicate name of award, awarding organizations and date)}
\begin{itemize}
  \item[] Reviewing awards: NeurIPS 2018 (top 7\%) and 2019 (``Top Reviewer''); ICML 2018 (``Outstanding Reviewer''), 2019 (``Best Reviewer''), 2020 (``Top Reviewer''), 2021 (``Best Reviewer'').
\end{itemize}

\nasubsection{
\subsection{Other Awards}
}
%
% In Computer Vision and Machine Learning venues, oral and spotlight presentations are reserved to the top papers chosen by the area chairs. Thus they are often considered prestigious.
% \begin{itemize}
%   \item \new{2021 --- Oral presentation (3.4\% acceptance rate), IEEE International Conference on Computer Vision (ICCV).}
%   \item 2019 --- Oral presentation (4.6\% acceptance rate), IEEE International Conference on Computer Vision (ICCV).
%   \item 2018 --- Oral presentation (2.1\% acceptance rate), IEEE/CVF Conference on Computer Vision and Pattern Recognition (CVPR).
%   \item 2016 --- Spotlight presentation (4.7\% acceptance rate), European Conference on Computer Vision (ECCV).
%   \item 2016 --- Oral presentation (3.9\% acceptance rate), IEEE/CVF Conference on Computer Vision and Pattern Recognition (CVPR).
% \end{itemize}

%%
%% 13. OTHER RELEVANT INFORMATION
%% Maximum One Page
\nasection{
  \ulsection{OTHER RELEVANT INFORMATION}

  Not applicable.
}


%%
%%%%%%%%%%%%%%%%%%%%%%%%%%%%%%%%%%%%%%%%%%%%%%%%%%%%%%%%%%%%%%%%%%%%%%%%
%       PUBLICATION RECORD
%%%%%%%%%%%%%%%%%%%%%%%%%%%%%%%%%%%%%%%%%%%%%%%%%%%%%%%%%%%%%%%%%%%%%%%%
\newpage
\UBCheader{Publication Record}\bigskip\rule{3.5in}{0mm}
\iftoggle{ubcformat}
{
\parbox[t]{2in}{{\bf Date:} \today}\parbox[t]{1in}{Initials: D.J.S.} \\
\noindent\parbox[t]{.5in}{\mbox{}}\parbox[t]{3.04in}{{\bf SURNAME:}
  Sutherland}\parbox[t]{3.2in}{{\bf FIRST NAME:} Danica}
\\\rule{3.54in}{0mm}\parbox[t]{3.2in}{{\bf MIDDLE NAME:} Jane}
\\[5mm]
\setcounter{section}{0}
\vspace{12pt}
}{}

% now included in header because it includes an \AtEveryBibitem
%\newcommand{\citeupdatedate}{ {{- "{:%B %Y}".format(cite_meta.last_cite_update) -}} }

\newcommand{\totalcites}{
  {%- set total_h = (cite_meta.total // 100) * 100 -%}
  {%- if cite_meta.total > total_h -%}over {% endif -%}
  {{- "{:,}".format(total_h) -}}
}
\newcommand{\totalcitesrec}{
  {%- set total_5y_h = (cite_meta.total_5y // 100) * 100 -%}
  {%- if cite_meta.total_5y > total_5y_h -%}over {% endif -%}
  {{- "{:,}".format(total_5y_h) -}}
}

\newcommand{\hindex}{ {{- cite_meta.h_index -}} }
\newcommand{\hindexrec}{ {{- cite_meta.h_index_5y -}} }

\newcommand{\itenindex}{ {{- cite_meta.i10_index -}} }
\newcommand{\itenindexrec}{ {{- cite_meta.i10_index_5y -}} }

\newcommand{\hfiveinfo}{
  {%- for venue in venues | top_h5_venues -%}
    {%- if loop.last and not loop.first %} and {% endif -%}
    {{ venue.short }} as \#{{ venue.h5_rank }}
    {%- if not loop.last %}, {% endif -%}
  {%- endfor -%}
}

\newcommand{\lotsoauths}{\parencite{
  {%- for paper in papers if paper.authors|length > 12 -%}
    {{- paper | bibtex_key(coauthors) -}}
    {%- if not loop.last -%},{%- endif -%}
  {%- endfor -%}
}}

\newcommand{\citethreshes}{
  {%- for thresh in cite_threshes|sort -%}
    {{- thresh -}}
    {% if not loop.last -%}, {% endif -%}
  {%- endfor -%}
}

\AtEveryBibitem{ {{- "" -}} %
  \csappto{blx@bbx@\thefield{entrytype}}{ {{- "" -}} % put at end of entry
    \rightsymb{ {{- "" -}} %
      \iffieldundef{usera}{\space}{\space\printfield{usera}}%
      {% for thresh in cite_threshes -%}
      \ifkeyword{cites {{- thresh -}} }{
        {%- if thresh|int() >= 100 -%} \textbf{ {%- endif -%}
        \citestext{ {{- "" -}} $\ge $ {{ "{:,}".format(thresh) }} {{- "" -}} }
        {%- if thresh|int() >= 100 -%} } {%- endif -%}
      }{}%
      {% endfor -%}
    }%
    \printfield{abstract}%
  }
}


\paragraph{Student authors.} Authors marked below in bold were my formal supervisees for the work that resulted in the paper in question.
% (Informal co-supervisees, discussed in section 8 \ref{sec:other-sup}, are not included.)

\paragraph{Author order.} Machine learning publications typically first list those primarily performing the bulk of the work (usually students and postdocs) in decreasing order of contribution,
followed by supervisors in increasing order of contribution and seniority.
Shared first authorship is denoted below with $\bibeqcon$.
In my few papers with a large number of co-authors -- \lotsoauths{} -- the full author list is suppressed for brevity,
but author order (outside of the first) is arbitrary, and none of these papers include my students as co-authors.

\paragraph{Publication venues.} The primary venues in machine learning work are conferences,
which conduct rigorous review processes with low acceptance rates (which are listed below where known).
NeurIPS and ICML are traditionally considered the most prestigious;
ICLR, AISTATS, and COLT are also primary venues in my subfield.
Although a deeply flawed metric, Google Scholar's h5-index (as of \citeupdatedate{}) lists \hfiveinfo{} among all publication venues in \emph{any} field.

\paragraph{Citation metrics.}
Citation metrics are from Google Scholar as of \citeupdatedate{}.
My work has been cited \totalcites{} times (last five years: \totalcitesrec{}),
with an $h$-index of \hindex{} (last five years: \hindexrec{}),
and \itenindex{} papers with at least ten citations (last five years: \itenindexrec{}).
Highly-cited papers are noted below (thresholds: \citethreshes).

% \paragraph{Citation stats.} \new{According to Google Scholar (visited April 25th, 2022)}:
% \vspace{-0.5em}
% \begin{itemize}
%   \setlength\itemsep{-0.5em}
%   \item \new{Total citations: 5027 (4367 since 2017)}
%   \item \new{h-index: 23 (21 since 2017)}
%   \item \new{i10-index: 33 (30 since 2017)}
% \end{itemize}

% \iftoggle{annualReview}{ Papers annotated as {\newpic}
%   are papers that were not included in the list last year that are fully
%   published, unless they are the under submitted work section---in these cases
%   they are new submissions. Papers marked as {\toappearpic} are ones that have
%   been accepted and have their camera ready versions submitted. Papers marked as
%     {\acceptedpic} are accepted but require revision for the camera ready version.
%   Papers marked as {\changedpic} are ones that were in any of the previous states
%   that had their states change within the last year. For example, {\changedpic}
%   in conference proceedings are in most cases submissions from last year that
%   have been accepted and published this year. }

\ulsection{REFEREED PUBLICATIONS}

\subsection{Journals}

\newrefcontext[labelprefix=J]
\printbibliography[heading=none,filter=journal]

\subsection{Conference Proceedings}

\newrefcontext[labelprefix=C]
\printbibliography[heading=none,filter=conf]

\subsection{Workshop Proceedings}
\remark{Workshops in machine learning play a similar role to conferences in most fields;
submissions are peer-reviewed but acceptance rates are high (and typically not published).
Most of these papers have (or will have) a corresponding conference version above.}

\newrefcontext[labelprefix=W]
\printbibliography[heading=none,filter=workshop]

\nasubsection{
\subsection{Other}

% \newrefcontext[labelprefix=O]
% \printbibliography[heading=none,filter=other]
}


\ulsection{NON-REFEREED PUBLICATIONS}

\nasubsection{
\subsection{Journals}
}

\nasubsection{
\subsection{Conferences}
\newrefcontext[labelprefix=CN]
\printbibliography[heading=none,filter=confnoreview]
}

\subsection{Other}
\newrefcontext[labelprefix=N]
\printbibliography[heading=none,filter=noreview]


\nasection{%
\ulsection{BOOKS}

\nasubsection{%
  \subsection{Authored}
}%

\nasubsection{%
  \subsection{Edited}
}

\nasubsection{%
  \subsection{Chapters}
}
}


\nasection{%
\ulsection{PATENTS}

% \newrefcontext[labelprefix=P]
% \printbibliography[heading=none,filter=patent]
}

\nasection{%
  \ulsection{SPECIAL COPYRIGHTS}

  Not applicable.
}%

\nasection{%
  \ulsection{ARTISTIC WORKS, PERFORMANCES, DESIGNS}

  Not applicable.
}%

\nasection{%
  \ulsection{OTHER WORKS}

  Not applicable.
}%

\nasubsection{
\ulsection{WORK SUBMITTED} % (including publisher and date of submission)}

% \newrefcontext[labelprefix=S]
% \printbibliography[heading=none,filter=submission]
}

\nasection{%
  \ulsection{WORK IN PROGRESS} % (including degree of completion)

  Not applicable.
}

\end{document}
